\section{Finite logarithmic hierarchies and reciprocal-logarithmic order}
\label{sec:logarithmic-scales}

A power gap and a logarithmic gap describe different orders of smallness.
For example, the correction $u/\log u$ has the same source power as $u$,
whereas $u^2\sqrt{-\log u}$ has a positive source-power gap and a
nonpolynomial logarithmic coefficient. These two situations, and a
leading product of logarithmic factors, require different normalizations. Exact
source-coordinate reductions, including repeated logarithmic charts, are
handled by the coordinate machinery of Section~\ref{sec:coordinates}.

\subsection{The finite coefficient hierarchy}

In the positive source coordinate $u\to0^+$, put
\begin{equation}
 L_1(u)=-\log u,\qquad L_{j+1}(u)=\log L_j(u),
 \qquad 1\le j<d.
 \label{eq:logarithmic-hierarchy}
\end{equation}
The depth $d$ is finite. All retained levels are positive and tend to
infinity sufficiently near the endpoint. A generalized logarithmic
monomial is
\begin{equation}
 c\prod_{j=1}^d L_j^{\alpha_j},\qquad
 c\in\R,\quad \alpha_j\in\R.
 \label{eq:generalized-log-monomial}
\end{equation}
Finite sums of these monomials form an algebra under addition and
multiplication. Powers in \eqref{eq:generalized-log-monomial} have positive
bases, including when an exponent is noninteger or negative. An arbitrary
real power of a sum is not silently included in this finite algebra.

\begin{proposition}[Closure under the source Euler derivative]
\label{prop:logarithmic-euler-closure}
For a differentiable coefficient $C(L_1,\ldots,L_d)$,
\begin{equation}
 u\frac{d}{du}C(L_1(u),\ldots,L_d(u))
 =\mathcal E C,
 \qquad
 \mathcal E=-\sum_{j=1}^d
 \frac{1}{L_1\cdots L_{j-1}}\frac{\partial}{\partial L_j},
 \label{eq:hierarchy-euler-operator}
\end{equation}
where the empty denominator is $1$. In particular, the finite generalized
logarithmic algebra is closed under every derivative used in the inverse
coefficient formula.
\end{proposition}
\begin{proof}
The chain rule gives $uL_1'(u)=-1$ and
$uL_j'(u)=-(L_1\cdots L_{j-1})^{-1}$. Each summand in
\eqref{eq:hierarchy-euler-operator} decreases finitely many monomial
exponents and creates no new logarithmic level.
\end{proof}

This is the needed closure for coefficients such as $L_1^{1/2}$:
its derivatives contain $L_1^{-1/2},L_1^{-3/2},\ldots$.
Restricting the coefficient algebra to positive logarithmic powers would
lose closure even in this simplest example. An expression containing
$\sqrt{\log u}$ instead has an eventually negative base on $u\to0^+$;
the real hierarchy does not admit it.

\subsection{Positive source-power gaps with logarithmic coefficients}

Consider
\begin{equation}
 F(u)=y_0+a u^p\left(1+\sum_{i=1}^m
 u^{\delta_i}C_i(L_1(u),\ldots,L_d(u))\right),
 \qquad a\ne0,\quad p\ne0,\quad \delta_i>0,
 \label{eq:generalized-log-forward}
\end{equation}
where every $C_i$ belongs to the finite algebra above. All exponents are
fixed real constants, and the coefficients and the sign of $a$ are fixed
on the parameter domain under consideration.
Set $z=((y-y_0)/a)^{1/p}>0$, so $z\to0^+$. For the observable $u^r$, the
multi-index coefficient of weight $w=\bk\cdot\bdelta$ is
\begin{equation}
 \frac{(-1)^{|\bk|}r}{p^{|\bk|}\prod_i k_i!}
 \prod_{j=1}^{|\bk|-1}(\mathcal E+r+w+pj)
 \prod_i C_i^{k_i},
 \label{eq:generalized-log-coefficient}
\end{equation}
evaluated at the logarithmic levels of $z$ and multiplied by $z^{r+w}$.
For $\bk=0$, the coefficient is $1$. Formula
\eqref{eq:generalized-log-coefficient} is the same Lagrange identity as
in the polynomial case; Proposition~\ref{prop:logarithmic-euler-closure}
supplies its enlarged coefficient algebra.

\begin{theorem}[Convergence and a conservative complete-tail bound]
\label{thm:generalized-log-convergence}
The inverse of \eqref{eq:generalized-log-forward} on its selected eventual
real branch has the convergent multi-index expansion
\eqref{eq:generalized-log-coefficient}. Choose integers $D_i\ge0$ such
that every monomial in $C_i$ satisfies
\[
 \prod_j L_j(u)^{\alpha_j}=\OO(\M(u)^{D_i}).
\]
For example, the ceiling of the largest sum of the positive monomial
exponents is sufficient. Let $\mathcal B$ be the finite minimal excluded
multi-index boundary of a positive source-weight cutoff, and put
\begin{equation}
 \beta=\min_{\bk\in\mathcal B}\bk\cdot\bdelta,
 \qquad D=\max_{\bk\in\mathcal B}\sum_i k_iD_i.
 \label{eq:generalized-log-boundary}
\end{equation}
Then the complete omitted expansion of $u^r$ is
$\OO(z^{r+\beta}\M(z)^D)$. At a finite source endpoint the original
inverse has $r=1$; at an infinite source endpoint it has $r=-1$,
with the selected source sign applied separately.
\end{theorem}
\begin{proof}
Write $u=zW$ near $W=1$. For each level, the quotient
$L_j(zW)/L_j(z)$ is analytic in $W$ and the finite reciprocal logarithmic
parameters near their limiting values. For instance, the first quotient
is $1-\log W/L_1(z)$; subsequent quotients are obtained by taking the
logarithm of the preceding quotient and dividing by the next level.
Each factor raised to a fixed real power is holomorphic near $1$.
Introduce a separate small parameter for every finite coefficient monomial
$z^{\delta_i}\prod_jL_j(z)^{\alpha_j}$. All tend to zero. The normalized
implicit equation has derivative $p\ne0$ in $W$ at its parameter origin,
so the analytic implicit function theorem proves convergence in a fixed
polydisc of these parameters.

Cauchy estimates, with the finite coefficient sums grouped by $i$, give a
geometric multi-index majorant in
$K_i z^{\delta_i}\M(z)^{D_i}$ for suitable constants $K_i$.
The omitted upward-closed index set is covered by the translates of its
finite minimal boundary. Sum the majorant over each translate. Its
denominators remain bounded for sufficiently small $z$, and the boundary
monomial is bounded by $z^\beta\M(z)^D$ using
\eqref{eq:generalized-log-boundary}. Multiplication by $z^r$ completes
the estimate. The estimate does not assume that the first boundary
coefficient is nonzero.
\end{proof}

The resulting logarithmic degree can be larger than the optimal one. The
finite boundary supplies a valid envelope for the entire omitted sum,
while the retained coefficients remain exact generalized expressions.
The cutoff uses the ordinary positive target coordinate: a source weight
$w$ contributes target power $(r+w)/|p|$. This convention preserves the
nonunit leading-power and infinite-endpoint scaling.

\begin{example}[Nonpolynomial logarithmic coefficients]
For $T=-\log y$, inversion at $0^+$ gives
\begin{align}
 (u+u^2\sqrt{-\log u})^{-1}(y)
 &=y-y^2\sqrt T+y^3(2T-\tfrac12)+\OO(y^4\M(y)^3),
 \label{eq:inverse-sqrt-log-coefficient}\\
 (u+u^2\log(-\log u))^{-1}(y)
 &=y-y^2\log T
 +y^3\left(2\log^2T-\frac{\log T}{T}\right)
 +\OO(y^4\M(y)^3).
 \label{eq:inverse-iterated-log-coefficient}
\end{align}
Here the superscript $-1$ denotes function inversion. The displayed
remainders are conservative complete-tail bounds.
The derivative term in each cubic coefficient is essential and follows
directly from \eqref{eq:hierarchy-euler-operator}.
\end{example}

\subsection{Reciprocal-logarithmic units}

Now consider a core of the form
\begin{equation}
 F_0(u)=y_0+a u^p H(-1/\log u),
 \qquad p\ne0,\quad H(0)=1,
 \label{eq:reciprocal-log-core}
\end{equation}
where $H$ is a rational function analytic at $0$ with real coefficients.
Let $z=((y-y_0)/a)^{1/p}>0$, $t=-1/\log z>0$, and write $u=zW$.
The normalized inverse equation is exactly
\begin{equation}
 W^p H\left(\frac{t}{1-t\log W}\right)=1.
 \label{eq:reciprocal-log-unit-equation}
\end{equation}
Its derivative in $W$ is $p$ at $(t,W)=(0,1)$, so its real solution
$W(t)$ is analytic at $0$. Iterating
\begin{equation}
 W\longmapsto
 H\left(\frac{t}{1-t\log W}\right)^{-1/p}
 \label{eq:reciprocal-log-fixed-point}
\end{equation}
from $W=1$ determines successively correct Taylor coefficients: dependence
on the previous $W$ carries a positive power of $t$. Exact finite Taylor
arithmetic determines the finite bracket and its first omitted orders.
An exclusive cutoff $h$ retains exactly the integer degrees $n<h$.
Its remainder has the form $\OO(t^\beta)$, where $\beta$ is the first
unretained nonzero degree, or any smaller degree beyond the retained terms
when no sharper order has been established.

For the original source inverse, the bracket is multiplied by the exact
prefactor $\sigma z$ at a finite endpoint or $\sigma/z$ at infinity.
An observable uses the corresponding power of $W$ and transports the
prefactor and remainder together. This is a logarithmic cutoff inside
that prefactor, not a cutoff in $y$ or $1/y$.

\begin{example}[A zero source-power gap]
At $u\to0^+$, the inverse of $u+u/\log u$ begins
\begin{equation}
 y\left(1+t+t^2+2t^3+\frac72t^4+\OO(t^5)\right),
 \qquad t=-\frac1{\log y}.
 \label{eq:reciprocal-log-inverse-example}
\end{equation}
Equivalently, its bracket is
$1-1/\log y+1/\log^2y-2/\log^3y+7/(2\log^4y)$.
For the same forward expression at $+\infty$, use $t=1/\log y$ and the
bracket $1-t+t^2-2t^3+7t^4/2$.

The expression $u+u/\log u-u/(1+\log u)$ has a cancelled first
reciprocal-logarithmic correction. Its forward unit is
$1+t^2/(1-t)$ in the source logarithmic coordinate, and its inverse bracket
begins $1-t^2+\cdots$. Coefficients are merged before detecting the leading
unit, so a cancelled coefficient does not create a fictitious retained term.
\end{example}

Higher source-power terms can accompany \eqref{eq:reciprocal-log-core}.
If each has a strictly positive power gap and an admitted logarithmic
coefficient envelope, their induced inverse correction is smaller than
the prefactor times every fixed power of $t$. For
$u(1+1/\log u)+u^2$, the logarithmic bracket therefore agrees with
\eqref{eq:reciprocal-log-inverse-example} at every finite logarithmic
order. The omitted source-power expression nevertheless belongs to a
distinct, smaller asymptotic sector. A residual computed solely from the
logarithmic bracket refers to the leading core. Equality of those brackets
does not assert that the full inverses are equal.

\subsection{A leading logarithmic monomial}

The same normalization extends to
\begin{equation}
 F_0(u)=y_0+a u^p\prod_{j=1}^d L_j(u)^{q_j},
 \qquad p\ne0,\quad q_j\in\R.
 \label{eq:leading-hierarchy-core}
\end{equation}
Let $Y=(y-y_0)/a>0$ and define
\begin{equation}
\begin{gathered}
 T=-\frac{\log Y}{p},\quad t=\frac1T,\quad
 B_1=T,\quad B_{j+1}=\log B_j,\\
 d_0=\frac1p\sum_jq_j\log B_j,
 \quad z=Y^{1/p}\prod_jB_j^{-q_j/p}.
\end{gathered}
 \label{eq:leading-hierarchy-normalization}
\end{equation}
All $B_j$ are positive sufficiently near the selected endpoint. Put
\begin{align}
 R_1(W)&=1+t(d_0-\log W),\label{eq:leading-level-first-ratio}\\
 R_{j+1}(W)&=1+B_{j+1}^{-1}\log R_j(W).
 \label{eq:leading-level-next-ratio}
\end{align}
Then $L_j(zW)/B_j=R_j(W)$ exactly, and the normalized inverse equation is
\begin{equation}
 W^p\prod_jR_j(W)^{q_j}=1.
 \label{eq:leading-hierarchy-unit-equation}
\end{equation}

\begin{theorem}[A convergent leading-hierarchy bracket]
\label{thm:leading-hierarchy-bracket}
Equation~\eqref{eq:leading-hierarchy-unit-equation} has a unique real
solution near $W=1$ for the selected asymptotic branch. Its expansion in
powers of $t$, retaining the lower-level coefficients exactly, converges
sufficiently near the endpoint. A truncation beginning its omitted degrees
at $\beta$ has bracket remainder
\begin{equation}
 \OO\!\left(t^\beta(1+|\log t|)^\beta\right).
 \label{eq:leading-hierarchy-tail}
\end{equation}
\end{theorem}
\begin{proof}
Treat $t$, $td_0$, and $B_2^{-1},\ldots,B_d^{-1}$ as independent small
parameters. The equation is analytic in these parameters and $W$ near
their origin and $W=1$, with derivative $p\ne0$ in $W$.
The analytic implicit function theorem gives a common polydisc and a
convergent series there. Since $d_0=\OO(\log T)$, all the parameters
indeed tend to zero on the actual asymptotic curve.

The total degree in the first two parameters equals the formal degree in
$t$ after substituting $td_0$. Cauchy estimates, uniform in the other
parameters on a smaller polydisc, bound the tail of those degrees by a
constant times $[t(1+|d_0|)]^\beta$ for sufficiently small $t$.
This proves \eqref{eq:leading-hierarchy-tail}. The fixed-point expression
$W=\prod_jR_j(W)^{-q_j/p}$ has positive formal valuation in its dependence
on the previous $W$, so finite exact iteration computes the same
coefficients. Real powers use the positive local factors near $1$.
\end{proof}

For $F(x)=x\log\log x$ at $+\infty$, let
$T=\log y$, $B=\log T$, and $t=1/T$. The reconstructed inverse starts
\begin{equation}
 F^{-1}(y)=\frac{y}{B}
 \left(1+\frac{\log B}{B}t
 +\OO\!\left(t^2(1+|\log t|)^2\right)\right).
 \label{eq:leading-iterated-log-example}
\end{equation}
Thus an iterated logarithm in the leading coefficient need not be replaced
by a constant. Products of levels and real powers of those levels are
covered by the same positive local normalization. Positive source-power
perturbations again belong to separate smaller sectors.

\subsection{Precision, exactness, and the boundary of the hierarchy}

The three normalizations above have different precision meanings. Positive
source-power gaps lead to a cutoff in the positive target power coordinate.
Reciprocal-logarithmic units and leading logarithmic monomials instead use
a cutoff in the small logarithmic coordinate inside an exact prefactor.
All these cutoffs are exclusive. Terms at a common source weight must be
combined before counting nonzero blocks; cancellation can remove a complete
block without removing its contribution to the index boundary used for the
tail estimate.

A finite run of zero coefficients does not establish termination. A
sufficient exactness argument reconstructs a positive candidate source from
the finite observable and verifies the original forward equation identically.
If no smaller sector has been omitted, local uniqueness then identifies
this candidate with the inverse germ. In particular, a polynomial source
unit does not by itself make a truncated noninteger power of that unit exact.

The normalized equations also give direct residual identities. Substituting
a finite unit into \eqref{eq:reciprocal-log-unit-equation} or
\eqref{eq:leading-hierarchy-unit-equation} determines its defect in the
leading-core equation. Such a defect is not a residual for a different
forward function containing unretained source-power sectors.

The reciprocal-logarithmic unit has a single ordinary Taylor bracket.
The other normalizations retain an enlarged coefficient algebra, so their
composition requires additional closure arguments. The convergence and
tail theorems give existence of constants and neighborhoods; they do not
supply explicit numerical values for them. An unknown forward remainder
likewise requires a separate transport estimate through the selected
logarithmic hierarchy.

These constructions form a finite hierarchy rather than a general
transseries field. Arbitrary real powers of coefficient sums, sums of
incomparable leading logarithmic monomials, and independent simultaneous
cutoffs at every logarithmic level require further hypotheses and additional
notions of precision. The formulas above establish closure under the Euler
derivatives used in inversion, convergence for the stated families, and
complete-tail estimates with the distinct precision meanings made explicit.
